\documentclass[11pt,a4paper]{article}
\usepackage[T1]{fontenc}
\usepackage{lmodern,amsmath,amssymb,amsthm,mathtools,booktabs,microtype}
\usepackage[margin=27mm]{geometry}
\usepackage[colorlinks=true,linkcolor=blue,citecolor=blue,urlcolor=blue]{hyperref}
\newtheorem{theorem}{Theorem}[section]
\newtheorem{lemma}[theorem]{Lemma}
\newtheorem{proposition}[theorem]{Proposition}
\newcommand{\Q}{\mathbb Q}\newcommand{\Z}{\mathbb Z}
\setlength{\emergencystretch}{2em}
\title{A quartic definition with 195 rational coordinates}
\author{}\date{16 November 2026}
\begin{document}\maketitle

\begin{abstract}
An affine change of the Lucas residual and its square permits a
Laurent elimination of an inverse-weight square. The resulting
system has 195 rational coordinates, including the input, and 175
quadratic equations. Its sum of integer squares has degree
four and content one. The transformation is a bijection on the full
rational zero sets, preserves the input, and retains the source's
arithmetic premises and sufficient parameter threshold. No additional
$abc$ hypothesis is introduced.
\end{abstract}

\section{Arithmetic scope and exact specification}

Let $C=149^N$, with $N\ge1$. Let $\mathcal S_{196}(C)$ be the
primitive quadratic system specified in~\cite{Primitive}, with its
input coordinate $X=x_1$. Its full coefficient table is
\begin{center}\small
\path{computation/quartic/primitive_basis196_v40/quadratic.json}.
\end{center}
The SHA-256 digest of that table is
\begin{center}\small\ttfamily
381ea54e77505a8d9b452b2afab42bd052beb0fb4dbd0804f63f2f609efbb275.
\end{center}
Its arithmetic assertion retains precisely the elliptic-group,
scalar contact-height, pole obstruction and prepared realization,
modular contact-height, and fixed-data simultaneous-prime premises
specified in~\cite{InverseVieta,Laurent}. These are the premises of the
following theorem. The finite polynomial verification below does not
independently certify them.

\begin{theorem}\label{thm:main}
Under the stated arithmetic premises, there is an integer $N_0$ such
that for every fixed $N\ge N_0$ there is a primitive integer quartic
$P_N$ in 195 rational coordinates, including $X$, for which
\[
 X\in\Z\quad\Longleftrightarrow\quad
 \exists y_2,\ldots,y_{195}\in\Q:\quad
 P_N(X,y_2,\ldots,y_{195})=0.
\]
The sufficient $N_0$ is exactly a sufficient threshold for
$\mathcal S_{196}$; it is ineffective. For every $N\ge1$, the
quadratic construction has 175 rows and its quartic has content one,
degree exactly four, and 2956 nonzero coordinate monomials. Its exact
coefficient height is
\[
 H(P_N)=\max\{126616122944180358034867798936,\;149^{2N+2}\}.
\]
The constant term in this maximum dominates for $1\le N\le5$;
the parameter term dominates for $N\ge6$. This crossover does not
specify the sufficient defining threshold $N_0$.
\end{theorem}

The complete polynomial is specified by the quadratic table and its
sum-of-squares expansion. The exact certificate records the affine
map, each Laurent coordinate exponent, the deleted coordinate, the
primitive row multiplier, and the old-to-new coordinate numbering.
These finite tables specify every coefficient, without an implicit
choice of equations or witnesses. They are distributed in
\begin{center}\small
\path{computation/quartic/affine_laurent195_v45/}.
\end{center}
The files \path{transformation.json}, \path{quadratic-family.json},
and \path{polynomial-family.json} give the maps, quadratic rows,
and complete quartic expansion, respectively.

\section{Elimination after an affine coordinate change}

We use the following elementary form of Laurent elimination.

\begin{lemma}\label{lem:laurent}
Let a polynomial system over $\Q$ have coordinates $x_1,\ldots,x_n$,
and let $u$ be a coordinate which is nonzero on every rational solution.
Retain that coordinate unchanged, with exponent zero, and let $x_d$
be a different coordinate.
Suppose an invertible affine change, fixing the input, is followed by
\[
 x_i=u^{a_i}y_i\quad(i\ne d),\qquad x_d=u^m,
\]
where all exponents are integers, the input exponent is zero, and
$x_d=u^m$ is a defining equation. Assume that the remaining equations,
after substitution and multiplication by nonzero Laurent monomials,
are polynomial equations of degree at most two. Assume also that these
new equations force $u\ne0$.
Then the new and old rational zero sets are in input-preserving
bijection. In particular they have the same rational input projection.
\end{lemma}

\begin{proof}
On the locus $u\ne0$ every displayed substitution is defined and
invertible. The eliminated coordinate has its unique specified value.
The affine change is invertible, and every row multiplier is nonzero.
Thus each direction of the coordinate map preserves and reflects all
equations. The nonzero assertion for the new system excludes any
additional boundary stratum. The input is fixed by both maps.
\end{proof}

For the source table, write
\[
 F=x_{50},\qquad F'=x_{51},\qquad
 S=x_{72},\qquad S'=x_{73}.
\]
Introduce centered coordinates $\widehat F,\widehat S,\widehat S'$ by
\begin{equation}\label{eq:affine}
 F=\widehat F+C^{-1},\qquad
 S=\widehat S+2\widehat F+C^{-1},\qquad
 S'=\widehat S'+F'.
\end{equation}
All remaining coordinates are fixed. This affine transformation is
triangular with determinant one for every nonzero $C$.
The source square equations
\[
 S=C(F^2+8F'^2),\qquad S'=CFF'
\]
therefore become
\[
 \widehat S=C(\widehat F^2+8F'^2),\qquad
 \widehat S'=C\widehat F F'.
\]
In particular the two combinations occurring in the detector factors
simplify exactly to
\[
 1-2CF+CS=C\widehat S,\qquad
 4CF'-4CS'=-4C\widehat S'.
\]
No square or auxiliary coordinate is added.

Set $w=x_{78}$ and $q=x_{77}$. The source contains the literal rows
\begin{equation}\label{eq:unit}
 qw=1,\qquad x_{79}=w^2.
\end{equation}
Apply the integer Laurent exponents listed in the certificate to the
centered coordinates, and eliminate $x_{79}$ using~\eqref{eq:unit}.
The first row of~\eqref{eq:unit} is retained unchanged. Consequently
$w\ne0$ on every rational solution of the new system, before any
Laurent division is made.

\begin{proposition}\label{prop:algebra}
For every nonzero rational $C$, the 175 certified rows in 195
coordinates are quadratic and have a rational zero set in
input-preserving bijection with that of $\mathcal S_{196}(C)$.
\end{proposition}

\begin{proof}
After~\eqref{eq:affine}, multiply each source row by the recorded
power of $C$ to clear its parameter denominators. Such a multiplier
is nonzero. For each remaining row, the certificate gives an integer
$e$ and a nonzero rational $c$ such that its Laurent substitution is
$cw^e$ times the specified target quadratic. The independent checker
verifies these identities as exact Laurent-polynomial identities.
It separately verifies that the deleted square equation becomes an
identity and that the inverse row $qw=1$ is retained literally.
The input exponent is zero. Lemma~\ref{lem:laurent} now applies.
\end{proof}

\section{Parameter scaling and coefficient height}

After the affine and Laurent transformation, make the fixed-parameter
coordinate change
\[
 y_i=C^{k_i}z_i,
\]
with the integer weights in the parameter-scaling certificate, and
divide each quadratic row by its recorded power of $C$. The input
weight is zero. Since $C=149^N\ne0$, this is another
input-preserving bijection on rational zero sets.
Every resulting row is an integer quadratic whose coefficients are
linear in $C$. The absolute value of every coefficient multiplying
$C$ in an individual row is at most $149$.

The final coefficient tables contain 2965 parameter-separated terms
and 2956 coordinate-monomial profiles. Each profile has degree at most
two in $C$. To certify that no coordinate monomial disappears at an
allowed parameter, the checker substitutes $C=149+t$ into the signed
profile: the constant term is strictly positive and the remaining
coefficients are nonnegative. Thus its sign is fixed for every
$C\ge149$.

The constant coefficient profile of greatest absolute value is
\[
 K=126616122944180358034867798936.
\]
The profile $22201C^2$ occurs. For $N=1,\ldots,5$, exact integer
substitution in every profile gives height $K$. For $C\ge149^6$, the
checker substitutes $C=149^6+t$ in each difference
$22201C^2\pm c_m(C)$, where $c_m$ is a coefficient profile; every
resulting coefficient is nonnegative. Hence the height is exactly
$22201C^2$ in this range. This proves the height formula in
Theorem~\ref{thm:main} for every $N\ge1$.

\section{The quartic and verification}

Denote the final certified integer quadratic rows by
$Q_1,\ldots,Q_{175}$. Define
\[
                 P_N=\sum_{j=1}^{175}Q_j^2.
\]
Over $\Q$, this polynomial vanishes precisely when all $Q_j$ vanish.
The input occurs only through its square in the retained input-square
row. Its squared leading term gives coefficient one on $X^4$, and
no other row contributes to that monomial. Therefore $P_N$ has degree
exactly four and integer content one for every allowed parameter.
Proposition~\ref{prop:algebra} and the source arithmetic assertion
prove Theorem~\ref{thm:main}.

The producer and independent checker use separate polynomial
implementations. The latter verifies the source digest, all 175 row
identities, the deleted identity, the nonzero guard, the fixed input,
the parameter-scaling identities, the complete sum-of-squares
expansion, and the coefficient-height and support certificates. These statements hold
symbolically in the fixed parameter; they are not tests at sampled
rational witnesses or numerical parameter values.

\begin{thebibliography}{9}
\bibitem{Primitive}
\emph{Primitive quadratic bases and small-coefficient quartic presentations},
\nolinkurl{papers/quartic/primitive_quadratic_bases_20261115_v40.tex}.
\bibitem{Laurent}
\emph{Quartics in 196 and 258 variables by elimination after monomial
rescaling},
\nolinkurl{papers/quartic/laurent_elimination_20261115_v38.tex}.
\bibitem{InverseVieta}
\emph{A quartic in 198 variables},
\nolinkurl{sources/h10_quartic198/arithmetic/inverse_vieta_198.tex}.
\end{thebibliography}
\end{document}
